\subsection{RGB to Luminance and Chrominance}

The RGB components may be transformed into three separate components,
luminance $Y$ and chrominance $Cb$ and $Cr$:
\begin{align}
Y  &= 0.299R + 0.587G + 0.114B\label{eq:y_luminance}\\
Cb &= -0.168736R - 0.331264G + 0.5B + 128\label{eq:cb_chrominance}\\
Cr &= 0.5R - 0.418688G - 0.081312B + 128\label{eq:cr_chrominance}
\end{align}

Here, $Y$ represents luminance, while $Cb$ and $Cr$ represent
chrominance information. The $Cb$ and $Cr$ components carry colour-difference
information associated with blue and red, respectively. They are separate
8-bit components in the usual 8-bit representation; they are not single-bit
colour indicators.
\vspace*{\baselineskip}
\subsection{Chrominance Components}

The $Cb$ component describes the blue-related colour difference and
$Cr$ describes the red-related colour difference. These components
retain colour information that is not required when only grayscale
intensity is required.

When chroma subsampling is used, the $Cb$ and $Cr$ components can be
stored at a lower spatial sampling rate than $Y$. For example, in 4:2:0
subsampling, each $2\times2$ group of pixels retains four $Y$ samples but
uses one $Cb$ sample and one $Cr$ sample. Thus, the chrominance components
remain separate components, while fewer chrominance samples are stored.
\vspace*{\baselineskip}
\subsection{Luminance Weighting}

The coefficients in the luminance equation are different because the
human visual system does not perceive the three RGB components with
 equal sensitivity:
\begin{align}
0.587&>0.299>0.114\label{eq:luminance_weights}
\end{align}

The green component consequently has the largest contribution to
luminance, followed by red and then blue.
\vspace*{\baselineskip}
\subsection{Information Reduction}

For the purpose of grayscale conversion, the colour pixel $(R,G,B)$ is
reduced to the single luminance value $Y$. The $Cb$ and $Cr$ components
carry the colour-difference information and are therefore not required
for the final grayscale intensity image.
